% TOP_PROBLEM: {"id":30007665,"problem_number":"LOCAL-30007665","title":"Durable anticorruption policy","category_id":21,"category":{"id":21,"name":"economics","display_name":"Economics","description":"Open economic theory statements, published research agendas, and proposed scoped specifications; question provenance and historical priority are qualified per record.","slug":"economics","order_index":21,"created_at":"2026-10-08T00:00:00Z"},"status":"research_question","external_url":null,"legacy_ids":[],"published":false,"statement":"In the explicitly specified setting: Which combinations of audits, sanctions, political competition, pay, and transparency reduce total corruption durably rather than displacing it into less observable forms? The mathematical statement above fixes the target, assumptions and scope of this catalogue formulation.","economics":{"original_id":154,"field":"Political economy and institutions","kind":"design","formalization_basis":"adapted_research_question","source_status":"explicit_open","priority_status":"earliest_found","priority_note":"Earliest directly inspected qualifying passage in this audit; earlier origins were not established. A review or research-agenda statement does not imply original priority.","earliest_verified_reference":{"authors":["Benjamin A. Olken","Rohini Pande"],"title":"Corruption in Developing Countries","year":2012,"url":"https://economics.mit.edu/research/publications/corruption-developing-countries","doi":"10.1146/annurev-economics-080511-110917","locator":"February 24, 2012 manuscript, section 4, printed pp. 35–37; conclusion, printed pp. 37–38","evidence_summary":"The review explicitly requests long-run anticorruption evidence allowing strategic adaptation and substitution among corruption forms."},"current_reference":{"authors":["Benjamin A. Olken","Rohini Pande"],"title":"Corruption in Developing Countries","year":2012,"url":"https://economics.mit.edu/research/publications/corruption-developing-countries","doi":"10.1146/annurev-economics-080511-110917","locator":"February 24, 2012 manuscript, section 4, printed pp. 35–37; conclusion, printed pp. 37–38","evidence_summary":"The review explicitly requests long-run anticorruption evidence allowing strategic adaptation and substitution among corruption forms."},"formalization_note":"The cited passage establishes a research direction, not this exact mathematical statement. The notation, population, policy menu, assumptions, and estimands are an original adaptation for this catalogue; no universal unsolved theorem or source priority is claimed."},"statusQualification":"Published question or research agenda verified at the cited locator. The mathematical specification is adapted for this catalogue and includes additional assumptions; source publication does not establish first-ever priority or certify this exact model remains unsolved. The cited passage establishes a research direction, not this exact mathematical statement. The notation, population, policy menu, assumptions, and estimands are an original adaptation for this catalogue; no universal unsolved theorem or source priority is claimed.","statusReviewedAt":"2026-10-08"}
% FIRST_POSTED: 2026-10-08
% ENTRY_KIND: statement_only
% !TeX program = lualatex
\documentclass[11pt]{article}
\usepackage{fontspec}
\usepackage[a4paper,margin=25mm]{geometry}
\usepackage{amsmath,amssymb,mathtools}
\usepackage{xurl}
\usepackage[hidelinks]{hyperref}
\newcommand{\sref}[2]{\hyperlink{source:#1:#2}{[S#2]}}
\setlength{\emergencystretch}{3em}
\begin{document}
% problemId:30007665
\section*{Durable anticorruption policy}\label{30007665}
\subsection{Definitions and mathematical statement}
Fix a public organization and a policy bundle \(p\) comprising audits, sanctions, transparency, competition, and personnel incentives. Let \(C^k_{jt}(p)\) be verified corruption through avenue \(k\) in unit \(j\) at date \(t\), and let \(\lambda_k\ge0\) be declared weights capturing losses rather than simply counts of detected offenses. Define total corruption harm \(C_{jt}(p)=\sum_k\lambda_kC^k_{jt}(p)\). The main estimand is \(\tau_C(T,p)=E[C_{jT}(p)-C_{jT}(0)]\), for a horizon permitting officials to learn and substitute across avenues. Estimate service output, enforcement cost, selection into office, and citizen welfare alongside \(\tau_C\). Audits may change detection probabilities, so measured offenses require an explicit measurement model or independent validation. A policy that lowers the targeted avenue while increasing another need not reduce total harm. Use credible bundle and follow-up variation to distinguish deterrence, adaptation, and personnel selection; describe interference across agencies. Determine which combinations remain effective after officials understand the rules. The source explicitly requests short- and long-run evidence under strategic adaptation. The avenue weights, outcome definition, and comparative policy objective are our formalization of that research agenda.

\par\textbf{Formulation and scope.} The cited passage establishes a research direction, not this exact mathematical statement. The notation, population, policy menu, assumptions, and estimands are an original adaptation for this catalogue; no universal unsolved theorem or source priority is claimed.

\par\textbf{Open-question provenance.} An explicit question or research agenda was verified at the source locator. This does not by itself certify that every later version remains unresolved \sref{30007665}{1}.

\par\textbf{Historical priority.} Earliest directly inspected qualifying passage in this audit; earlier origins were not established. A review or research-agenda statement does not imply original priority.

\subsection{Short English statement}
In the explicitly specified setting: Which combinations of audits, sanctions, political competition, pay, and transparency reduce total corruption durably rather than displacing it into less observable forms? The mathematical statement above fixes the target, assumptions and scope of this catalogue formulation.

\subsection{Sources}
\begin{itemize}\raggedright
\item[\textnormal{[S1]}]\hypertarget{source:30007665:1}{} Benjamin A. Olken, Rohini Pande. 2012. Corruption in Developing Countries. February 24, 2012 manuscript, section 4, printed pp. 35–37; conclusion, printed pp. 37–38.
\url{https://economics.mit.edu/research/publications/corruption-developing-countries}.
DOI: 10.1146/annurev-economics-080511-110917.
{\small\textit{Earliest verified question/agenda reference; historical priority qualified below. The review explicitly requests long-run anticorruption evidence allowing strategic adaptation and substitution among corruption forms.}}
\end{itemize}
\end{document}
